\documentclass[11pt]{article}

\usepackage[margin=1in]{geometry}
\usepackage{amsmath,amssymb}
\usepackage{booktabs}
\usepackage{enumitem}
\usepackage[hidelinks]{hyperref}

\setlist[itemize]{leftmargin=1.4em,topsep=2pt,itemsep=2pt}
\setlist[enumerate]{leftmargin=1.6em,topsep=2pt,itemsep=3pt}

\newcommand{\Iz}{I_{0}}
\newcommand{\phz}{\varphi_{0}}
\newcommand{\Vp}{V_{\mathrm{M+}}}
\newcommand{\Vm}{V_{\mathrm{M-}}}
\newcommand{\Resp}{\mathcal{R}}

\title{\vspace{-1.5em}Experimental Proposal: Measurement of the Total Optical
Intensity and Static Phase Bias of the AWA TFLN Electro-Optic THz Sensor}
\author{BeamNetUS / SULI THz Diagnostics}
\date{\today}

\begin{document}
\maketitle
\vspace{-1em}

\section{Purpose}

The experiment at the Argonne Wakefield Accelerator performed previously yielded
results that are not able to be interpreted fully without further parameters that can
only be obtained experimentally. Specifically, the combined laser intensity from both
branches of the TFLN sensor and the internal phase bias are unknown currently, and
these parameters are required to perform an accurate analysis of the electric field of
the coherent transition radiation incident on the sensor. The purpose of this
experiment is to determine these parameters by recreating the necessary components of
the previously performed experiment and performing measurements.

The two unknown parameters are the total optical intensity reaching the balanced
photodiode pair, $\Iz$, and the static phase bias of the Mach--Zehnder interferometer,
$\phz$. Both are properties of the laser, fiber, and device together, not of the
electron beam. The slope that converts a recorded scope voltage into an electro-optic
phase shift scales as $\Iz \sin(2\phz)$, so an uncertainty in either parameter carries
through directly into the recovered field. Neither was recorded during the previous
runs, which is why a dedicated measurement is needed.

These parameters cannot be read from a specification. The TFLN device is a poled
Mach--Zehnder interferometer: during fabrication one arm is domain-inverted by applying
a field above the lithium-niobate coercive field of about 22 kV/mm, after which the
poling electrodes are removed. The finished sensor has no bias electrodes and is
operated with no applied bias. The static phase bias is therefore set only by passive
effects, including residual fabrication imbalance, the fiber and polarization state, and
thermal drift. It must be measured in place, and it may change between runs.

\section{Measurement Principle}

When no THz field is present, the two complementary outputs of the interferometer carry
optical powers
\begin{equation}
  P_{1} = \Iz \cos^{2}\!\phz, \qquad P_{2} = \Iz \sin^{2}\!\phz,
\end{equation}
so that
\begin{equation}
  P_{1}+P_{2} = \Iz, \qquad P_{1}-P_{2} = \Iz \cos(2\phz).
\end{equation}
Diode 1 feeds the MONITOR$+$ output and diode 2 feeds the MONITOR$-$ output, while the
RF output reports the difference between the two. The DC-coupled MONITOR voltages are
\begin{equation}
  \Vp = \Resp\,G_{\mathrm{M}}\,P_{1}, \qquad \Vm = \Resp\,G_{\mathrm{M}}\,P_{2},
\end{equation}
where $\Resp$ is the photodiode responsivity at 1550 nm (about 1 A/W) and
$G_{\mathrm{M}}$ is the MONITOR transimpedance gain (10 kV/A into a high-impedance load,
or 1.55 V/mW into 50 $\Omega$).

The two MONITOR readings determine both unknowns. The total intensity is
\begin{equation}
  \Iz = \frac{\Vp + \Vm}{\Resp\,G_{\mathrm{M}}}, \label{eq:I0}
\end{equation}
and the static phase bias follows from the balance of the two arms,
\begin{equation}
  \cos(2\phz) = \frac{\Vp - \Vm}{\Vp + \Vm}, \qquad
  \phz = \tfrac{1}{2}\arccos\!\left(\frac{\Vp - \Vm}{\Vp + \Vm}\right). \label{eq:phi0}
\end{equation}
Equation \eqref{eq:phi0} is a ratio, so the bias angle does not depend on the
responsivity, the gain, or the intensity. Equal MONITOR readings mean the device sits
at quadrature ($\phz = 45^\circ$), where the response to a small phase shift is largest;
a large imbalance means it sits near a cosine extremum, where the response is small. The
bias factor that sets the slope is
\begin{equation}
  s \equiv \sin(2\phz) = \sqrt{1-\left(\frac{\Vp-\Vm}{\Vp+\Vm}\right)^{2}}
    = \frac{2\sqrt{\Vp\,\Vm}}{\Vp+\Vm}. \label{eq:s}
\end{equation}

The quantity used in the downstream analysis is the BPD calibration slope, the RF output
voltage per radian of electro-optic phase,
\begin{equation}
  K_{\mathrm{BPD}} = 2\,G_{\mathrm{eff}}\,\Resp\,\Iz\,\sin(2\phz),
\end{equation}
where $G_{\mathrm{eff}}$ is the RF-output effective gain ($G_{\mathrm{switch}}/2$ into
50 $\Omega$). Substituting \eqref{eq:I0} and \eqref{eq:s}, the responsivity and the
intensity cancel, and the calibration reduces to
\begin{equation}
  K_{\mathrm{BPD}} = 4\,\frac{G_{\mathrm{eff}}}{G_{\mathrm{M}}}\,\sqrt{\Vp\,\Vm}.
  \label{eq:Kbox}
\end{equation}
This is the result that motivates the measurement. The full voltage-to-phase calibration
is fixed by the two recorded MONITOR voltages and the known ratio of the RF and MONITOR
gains. The responsivity and the absolute intensity, which were the two largest sources
of uncertainty in the prior analysis, drop out of $K_{\mathrm{BPD}}$ as long as the
MONITOR and RF outputs are recorded at the same time. The gain-switch position is the
one parameter that still must be recorded by hand.

\section{Equipment}

The laser setup at the Argonne Wakefield Accelerator is a fiber-coupled erbium-doped
laser, operating at 81.25 MHz with 100 mW average power delivered in 70-fs pulses. The
laser operates at 1550 nm.

Partow Technologies developed a thin-film lithium niobate (TFLN) electro-optic sensor
that uses a Mach--Zehnder interferometer to induce an electric-field-dependent phase
modulation on incident laser light. The device is fiber-coupled, has two complementary
outputs, and has a measured optical insertion loss of $-13$ dB at 1550 nm.

Detection of the MZI outputs is performed by a Thorlabs PDB450C-AC balanced photodiode.
This device measures the intensity of two incident fiber-coupled sources and reports the
difference between the two through a balanced RF output. It also provides individual
monitor outputs for the two separate diode channels. The RF output is AC-coupled with a
switched transimpedance gain of $10^{3}$ to $10^{7}$ V/A. The MONITOR$+$ and MONITOR$-$
outputs are DC-coupled at 10 kV/A with a bandwidth of DC to 1 MHz.

The acquisition system is an oscilloscope with at least three channels. The RF output is
placed on a 50 $\Omega$-terminated channel. The MONITOR outputs are placed on
high-impedance (1 M$\Omega$) channels or read with a precision DVM, which is sufficient
for their DC-to-1-MHz bandwidth and avoids the factor-of-two attenuation of a 50 $\Omega$
input. The fiber link uses the same type and length of polarization-maintaining fiber as
the AWA setup, nominally 25 m, to reproduce the optical loss and polarization
conditions. A calibrated 1550 nm optical power meter is used for the intensity
cross-check.

\section{Experimental Arrangements}

The optical path outlined in ``Implementation of a 1550-nm Laser System for Beam
Characterization at the Argonne Wakefield Accelerator'' should be recreated as
faithfully as possible, without installation at the accelerator. The laser should be
coupled to the TFLN sensor with the same type and quantity of fiber cabling. The two MZI
outputs should be connected to the two inputs of the PDB450C-AC. All three RF outputs
should be monitored by the oscilloscope: the balanced output, MONITOR$+$, and
MONITOR$-$. The laser should be powered on in the same operational state as before, the
gain-switch position should be recorded, and results should be recorded.

\subsection{Primary measurement}
With no THz field present, the MONITOR$+$ voltage, the MONITOR$-$ voltage, and the RF
baseline voltage should be recorded at the same time. The intensity $\Iz$ is then
computed from \eqref{eq:I0}, the static phase bias $\phz$ from \eqref{eq:phi0}, the bias
factor $s$ from \eqref{eq:s}, and the calibration $K_{\mathrm{BPD}}$ from
\eqref{eq:Kbox}. As an internal check, the recorded RF baseline should satisfy
$V_{\mathrm{RF,0}} = (G_{\mathrm{eff}}/G_{\mathrm{M}})(\Vp-\Vm)$, which confirms the gain
ratio and the channel assignments.

\subsection{Intensity cross-check}
The total optical power at the device input should be measured with the calibrated power
meter and propagated through the $-13$ dB insertion loss to predict
$\Iz \approx 0.05 \times P_{\mathrm{in}}$. This value should be compared with the
intensity obtained from the MONITOR readings. Agreement bounds the systematic error on
the intensity scale, and a disagreement points to its cause, such as an incorrect
insertion loss, a nonlinear MONITOR response, or a responsivity offset.

\subsection{Auxiliary measurements}
\begin{itemize}
  \item Responsivity cross-calibration. The laser should be connected directly to one
        BPD input through a calibrated attenuator, and the MONITOR reading compared with
        the power meter, to fix the responsivity at 1550 nm to about 5 percent. This
        improves the intensity from \eqref{eq:I0} but is not required for the
        calibration in \eqref{eq:Kbox}.
  \item Linearity and saturation check. The launched power should be varied to confirm
        that the MONITOR voltages scale linearly. The MONITOR stage saturates near 1 mW
        of optical power per diode, as discussed below.
  \item Drift characterization. The MONITOR voltages should be logged over tens of
        minutes to hours to measure how quickly the phase bias wanders. This sets how
        often the bias must be re-sampled during a beam run.
  \item Bypass baseline. The full chain may be recorded with the sensor removed, with
        the laser sent into the BPD through the rest of the optics, to give an upper
        bound on the intensity. The ratio of this reading to the in-device reading gives
        a measured insertion loss to compare against the $-13$ dB value.
\end{itemize}

\section{Analysis}

For each acquisition, the intensity $\Iz$, the phase bias $\phz$, the bias factor
$s = \sin(2\phz)$, and the calibration $K_{\mathrm{BPD}}$ should be reported with
uncertainties propagated from the measured voltages and the datasheet gains. From
\eqref{eq:Kbox} the relative uncertainty in the calibration is
\begin{equation}
  \frac{\delta K_{\mathrm{BPD}}}{K_{\mathrm{BPD}}}
   = \sqrt{\left(\frac{\delta(G_{\mathrm{eff}}/G_{\mathrm{M}})}{G_{\mathrm{eff}}/G_{\mathrm{M}}}\right)^{2}
   + \tfrac{1}{4}\left(\frac{\delta\Vp}{\Vp}\right)^{2}
   + \tfrac{1}{4}\left(\frac{\delta\Vm}{\Vm}\right)^{2}},
\end{equation}
so few-percent voltage readings give a few-percent calibration. The intensity from the
MONITOR readings and from the power meter should agree within their combined error, and
the spread between them is reported as the systematic uncertainty on the intensity
scale.

\section{Feasibility}

Prior records indicate about 24 mW at the fiber-link exit. Through the $-13$ dB device,
about 1.2 mW reaches the BPD, or roughly 0.6 mW per diode. The expected MONITOR levels
are about 6 V into a high-impedance load, or about 0.9 V into 50 $\Omega$, per arm, which
are comfortably readable.

The main feasibility risk is saturation. At 0.6 mW per diode the MONITOR stage is at
about 60 percent of its 1 mW saturation point. If the linearity check fails, the launched
power should be attenuated, or the coupling reduced, until both arms are in the linear
regime. Because the phase bias in \eqref{eq:phi0} is a ratio and the calibration in
\eqref{eq:Kbox} uses the geometric mean of the two MONITOR voltages, moderate
attenuation does not bias the result as long as both arms stay linear.

A well-fabricated fiber-coupled interferometer is usually delivered near quadrature, so
the two MONITOR readings are expected to be similar and the bias factor close to one. The
measurement will confirm whether the AWA device sits near maximum sensitivity or has
drifted toward an extremum.

\section{Limitations and Implications}

The phase bias is passive and temperature-dependent, so the value measured on the bench
need not equal the value during any past AWA run. The bench experiment establishes the
method, the baseline values, and the drift rate, rather than a retroactive phase bias for
archived data.

If the intensity proves reproducible, set by the laser power and the fiber coupling, then
for an archived run the recorded RF baseline and the assumed gain give that run's bias
through $\cos(2\phz) = V_{\mathrm{RF,0}}/(G_{\mathrm{eff}}\Resp\Iz)$. This is offered as a
bounded estimate, with reproducibility verified by repeating the primary measurement
across re-couplings.

The decisive outcome is a standing acquisition requirement for all future beam
campaigns: record MONITOR$+$, MONITOR$-$, and the RF output at the same time, along with
the gain-switch position. Equation \eqref{eq:Kbox} then gives an absolute,
responsivity-independent field calibration for every shot, which closes the largest open
uncertainty in the AWA electro-optic sampling analysis.

\section{Summary}

By recreating the AWA optical chain on the bench and recording the two PDB450C-AC MONITOR
outputs alongside the RF output, this experiment measures the total optical intensity and
the static phase bias directly, checks the intensity against a power-meter measurement,
and shows that the voltage-to-phase calibration is fixed by the two MONITOR voltages and
the known detector gain ratio. The result both constrains the existing dataset and
defines the acquisition protocol that makes future electro-optic sampling data
absolutely calibratable.

\end{document}
